\documentclass{article}
\pagestyle{empty}
\usepackage{amsmath}
\usepackage{geometry}
\geometry{a4paper, margin=1.5cm}
%\title{Activité : Nourrir les dauphins}
\date{}
\begin{document}

\begin{center}
\section*{Activité 1: Étude du saut d'un dauphin}
\end{center}

Dans un bassin, un dauphin nage tranquillement et aperçoit, à l’instant $t = 0$ secondes, un poisson lancé par son soigneur. Il saute, l’attrape 1 seconde plus tard à une hauteur de 1,5 mètre et replonge dans l’eau au bout de 6 secondes.

La trajectoire de son saut est modélisée par la courbe d'une fonction $f$ définie par :
$$
f(t) = a (t - t_1)(t - t_2),
$$
où $a$, $t_1$ et $t_2$ désignent des réels et $a \neq 0$.

Le niveau du plan d’eau du bassin correspond à l’altitude 0 mètre et $f(t)$ est la hauteur, en mètre, atteinte par le dauphin au bout d’une durée $t$, en secondes.

\begin{enumerate}
    \item En observant la forme de l'équation $f(t) = a (t - t_1)(t - t_2)$, montrez que résoudre $f(t) = 0$ implique que soit $t - t_1 = 0$, soit $t - t_2 = 0$. Quelles valeurs cela donne-t-il pour $t_1$ et $t_2$ dans le contexte du saut du dauphin ?

    \item Pour déterminer le coefficient $a$, utilisez le fait que le dauphin atteint une hauteur de 1,5 mètre après 1 seconde.

    \item Développez cette expression.
\end{enumerate}

\hrule

\begin{center}
\section*{Activité 1: Étude du saut d'un dauphin}
\end{center}

Dans un bassin, un dauphin nage tranquillement et aperçoit, à l’instant $t = 0$ secondes, un poisson lancé par son soigneur. Il saute, l’attrape 1 seconde plus tard à une hauteur de 1,5 mètre et replonge dans l’eau au bout de 6 secondes.

La trajectoire de son saut est modélisée par la courbe d'une fonction $f$ définie par :
$$
f(t) = a (t - t_1)(t - t_2),
$$
où $a$, $t_1$ et $t_2$ désignent des réels et $a \neq 0$.

Le niveau du plan d’eau du bassin correspond à l’altitude 0 mètre et $f(t)$ est la hauteur, en mètre, atteinte par le dauphin au bout d’une durée $t$, en secondes.

\begin{enumerate}
    \item En observant la forme de l'équation $f(t) = a (t - t_1)(t - t_2)$, montrez que résoudre $f(t) = 0$ implique que soit $t - t_1 = 0$, soit $t - t_2 = 0$. Quelles valeurs cela donne-t-il pour $t_1$ et $t_2$ dans le contexte du saut du dauphin ?

    \item Pour déterminer le coefficient $a$, utilisez le fait que le dauphin atteint une hauteur de 1,5 mètre après 1 seconde.

    \item Développez cette expression.
\end{enumerate}


\hrule

\begin{center}
\section*{Activité 1: Étude du saut d'un dauphin}
\end{center}

Dans un bassin, un dauphin nage tranquillement et aperçoit, à l’instant $t = 0$ secondes, un poisson lancé par son soigneur. Il saute, l’attrape 1 seconde plus tard à une hauteur de 1,5 mètre et replonge dans l’eau au bout de 6 secondes.

La trajectoire de son saut est modélisée par la courbe d'une fonction $f$ définie par :
$$
f(t) = a (t - t_1)(t - t_2),
$$
où $a$, $t_1$ et $t_2$ désignent des réels et $a \neq 0$.

Le niveau du plan d’eau du bassin correspond à l’altitude 0 mètre et $f(t)$ est la hauteur, en mètre, atteinte par le dauphin au bout d’une durée $t$, en secondes.

\begin{enumerate}
    \item En observant la forme de l'équation $f(t) = a (t - t_1)(t - t_2)$, montrez que résoudre $f(t) = 0$ implique que soit $t - t_1 = 0$, soit $t - t_2 = 0$. Quelles valeurs cela donne-t-il pour $t_1$ et $t_2$ dans le contexte du saut du dauphin ?

    \item Pour déterminer le coefficient $a$, utilisez le fait que le dauphin atteint une hauteur de 1,5 mètre après 1 seconde.

    \item Développez cette expression.
\end{enumerate}


\end{document}
